\section{Task Definitions and Success Criteria}
\label{app:tasks}

We provide additional details for the simulation and real-world experiments. All tasks require the robot to identify and interact with task-relevant functional parts, rather than relying only on object-level category recognition. We use the same task initialization protocol, maximum execution horizon, and success criteria for all compared methods.

\paragraph{Simulation tasks.}
\begin{itemize}
\item \textbf{Hanging Mug.} The robot grasps the rim or another stable graspable region of a mug and hangs the mug by placing its handle onto a target rack. This task requires identifying both a reliable grasping region and the mug handle used for hanging. A trial is successful if the mug is stably hung on the target rack without falling.

\item \textbf{Beat Block with Hammer.} The robot grasps the handle of a hammer and uses the hammer head to hit a target block. The task requires distinguishing the hammer handle from the hammer head and generating effective contact with the hammer head. A trial is successful if the hammer head contacts the block and moves it to the target state.

\item \textbf{Open Microwave.} The robot opens a microwave door by interacting with the door handle or another valid interactive region. The task requires localizing the articulated interaction region and applying motion in the correct direction. A trial is successful if the microwave door is opened beyond the predefined angle threshold.

\item \textbf{Put Object into Cabinet.} The robot first interacts with the cabinet handle to open the cabinet, and then places the target object into the cabinet or target storage region. This task requires identifying the cabinet handle as the functional interaction region and completing the subsequent placement. A trial is successful if the cabinet is opened and the object is placed inside the target cabinet region without falling.
\end{itemize}

\begin{figure*}[t]
    \centering
    \includegraphics[width=0.9\linewidth]{Images/simulation_experiment.png}
    \caption{Simulation task examples used in our evaluation. These tasks require localizing functional object parts, such as mug handles, hammer heads, articulated handles, and placement regions, to complete object-centric manipulation.}
    \label{fig:simulation_experiment}
\end{figure*}


\paragraph{Real-world tasks.}
\begin{itemize}
\item \textbf{Grasp Mug.} The robot grasps a mug from the table and lifts it stably. This task requires identifying a reliable graspable region, such as the mug handle, rim, or body. A trial is successful if the robot lifts the mug above a height threshold without slipping or dropping it.

\item \textbf{Beat Cube.} The robot grasps a hammer and uses the hammer head to hit a target cube. This task requires distinguishing the hammer handle from the hammer head and generating contact with the hammer head rather than the handle or side surface. A trial is successful if the hammer head contacts the cube and causes visible displacement or reaches the target state.

\item \textbf{Stir Mug.} The robot grasps a stirring tool and inserts its tip into a mug to perform a stirring motion. This task requires identifying the graspable region and functional tip of the tool, as well as the opening of the mug. A trial is successful if the tool tip enters the mug and completes the stirring motion without losing control or colliding severely with the mug.

\item \textbf{Pour Water.} The robot grasps a container and pours its contents into a target cup or target region. This task requires identifying a stable grasp region, the opening or pouring direction of the container, and the target container location. A trial is successful if the container is tilted toward the target and the contents enter the target cup or region.
\end{itemize}

\section{Real-World Setup and Point-Cloud Processing}
\label{app:real_world_setup}

For the real-world experiments, we use a bimanual robot platform equipped with two calibrated ZED 2i RGB-D cameras. Figure~\ref{fig:experiment_setup} shows the robot setup and the object instances used in our real-world evaluation. At each policy step, we reconstruct the scene point cloud from calibrated RGB-D observations and extract object-level point clouds for the task-relevant objects.

\begin{figure*}[t]
    \centering
    \includegraphics[width=0.9\linewidth]{Images/experiment_setup}
    \caption{Real-world experimental setup. The bimanual robot platform is equipped with two calibrated ZED 2i RGB-D cameras, and the evaluated objects include mugs, hammers, stirring tools, and containers used in the real-world manipulation tasks.}
    \label{fig:experiment_setup}
\end{figure*}

To obtain object point clouds, we use text-prompted real-time tracking and segmentation with SAM2~\cite{ravi2024sam2}. For each target object, a text prompt specifies the object of interest in the RGB observations. SAM2 tracks the corresponding object masks over time and across camera views. The segmented multi-view RGB-D observations are then back-projected into 3D using camera intrinsics and extrinsics, producing object-specific point clouds in the world coordinate frame. These object point clouds are used both as raw object observations for point-wise baselines and as the support condition for our semantic field. All compared methods share the same RGB-D observations, calibration, segmentation masks, and point-cloud extraction pipeline.

\begin{figure*}[t]
    \centering
    \includegraphics[width=0.9\linewidth]{Images/object_segmentation}
    \caption{Object point-cloud extraction pipeline. Text prompts specify task-relevant objects in multi-view RGB observations; SAM2 tracks and segments the objects over time; the masked RGB-D observations are back-projected with camera calibration to obtain object-specific point clouds for policy input and semantic-field conditioning.}
    \label{fig:object_segmentation}
\end{figure*}

\section{Representation and Policy Implementation Details}
\label{app:representation_policy_details}

\paragraph{Semantic field training.}
We train the semantic field from part-annotated PartNext object models \cite{wang2026partnext}. For each object family, all training instances use a consistent part label set, such as handle/body for mugs or head/handle for hammers. The Utonia backbone is initialized from the public checkpoint and kept frozen; the trainable components are the lightweight feature adapter, tri-plane feature fusion module, and query decoder. During training, support points are sampled from the object surface to condition the field, and query points are sampled from labeled surface locations to provide part supervision. For augmentation stability, we generate two augmented support conditions from the same object using independent SO(3) rotations and Gaussian coordinate jitter, while supervising corresponding query points across the two views. Table~\ref{tab:semantic_field_hparams} summarizes the main representation training and architecture settings.

\begin{table}[h]
\centering
\small
\caption{Semantic field training and architecture settings used in our experiments.}
\label{tab:semantic_field_hparams}
\begin{tabular}{@{}p{0.32\linewidth}p{0.62\linewidth}@{}}
\toprule
Component & Setting \\
\midrule
Dataset & PartNext meshes with part annotations \\
Trainable components & Feature adapter, tri-plane fusion, and query decoder \\
Backbone & Frozen Utonia checkpoint \\
Support points & 5000 surface points per object \\
Query points & 2048 labeled surface points per object \\
Coordinate augmentation & SO(3) rotation and Gaussian jitter, $\sigma=0.005$ \\
Tri-plane cache & Resolution $64$, padding $0.05$ \\
Adapter / branch dimension & 256 / 256 \\
Semantic embedding dimension & 128, L2-normalized \\
Query decoder & Two LayerNorm--GELU MLP layers with embedding and logit heads \\
Optimizer & AdamW, learning rate $10^{-3}$, weight decay $10^{-4}$ \\
Batch size / epochs & 6 / 4000 \\
Loss weights & $\lambda_{part}=0.5$, $\lambda_{align}=0.2$, $\lambda_{stab}=0.1$ \\
Mixed precision & Enabled \\
\bottomrule
\end{tabular}
\end{table}

\paragraph{Semantic point-cloud export.}
After semantic field training, the field is frozen and used only as a representation module. For each observed object point cloud, we first remove invalid zero points and use the remaining object points as the support condition. We then sample a fixed number of object query locations and evaluate the embedding output of the field. In the default policy preprocessing, we export $256$ query points per semantic object, and each exported point stores its world-frame coordinate concatenated with a 128-dimensional semantic embedding, yielding an object-level semantic point cloud of shape $256\times(3+128)$. Part logits are not passed to the policy; they are used only during representation training and optional visualization.

\paragraph{Policy integration.}
We use DP3 as the policy backbone and keep its diffusion objective unchanged. In the observation encoder the scene point cloud and each exported semantic point cloud are treated as separate point-cloud modalities. The scene point cloud provides global geometry, while each semantic point cloud provides XYZ coordinates
concatenated with semantic embeddings for a target object. The encoded point-cloud features are concatenated with the robot proprioceptive feature to form the conditioning vector for the diffusion policy. Table~\ref{tab:policy_hparams} summarizes the downstream policy and semantic point-cloud preprocessing settings.



\begin{table}[h]
\centering
\small
\caption{Downstream policy and semantic point-cloud preprocessing settings.}
\label{tab:policy_hparams}
\begin{tabular}{@{}ll@{}}
\toprule
Component & Setting \\
\midrule
Policy backbone & DP3 point-cloud diffusion policy \\
Scene point cloud & 1024 points \\
Semantic point clouds & Up to two target objects, 256 points each \\
Semantic point dimension & $3+128$ channels \\
Observation horizon & 3 steps \\
Action horizon & 6 steps, with total horizon 8 \\
Action / proprioception dimension & 14 / 14 \\
Point-cloud encoder & PointNet encoder, output dimension 128 per modality \\
Diffusion scheduler & DDIM, 100 training steps and 10 inference steps \\
Policy optimizer & AdamW, learning rate $10^{-4}$, weight decay $10^{-6}$ \\
Batch size / epochs & 256 / 3000 \\
\bottomrule
\end{tabular}
\end{table}

\paragraph{Hardware.}
All representation and policy models are trained on a desktop workstation with a single NVIDIA RTX 4090 GPU and Ubuntu 22.04. Real-robot deployment and online policy inference are run on a laptop with an NVIDIA RTX 4070 GPU and Ubuntu 22.04. The semantic field is frozen during deployment, so online computation only requires object point-cloud extraction, field querying for the semantic point clouds, and DP3 policy inference.

\paragraph{Runtime and asynchronous execution.}
During real-world deployment, the perception and policy pipeline includes object point-cloud construction, SAM2-based object segmentation, semantic/observation encoding, and DP3 policy inference. Table~\ref{tab:runtime_breakdown} reports the measured latency of the main computation stages on the RTX 4070 laptop used for deployment. The end-to-end wall-clock time of a full inference update can be higher than the sum of these stages due to camera synchronization, robot I/O, data transfer, and action post-processing.

\begin{table}[h]
\centering
\small
\caption{Runtime breakdown of the main computation stages during real-world deployment.}
\label{tab:runtime_breakdown}
\begin{tabular}{@{}lc@{}}
\toprule
Stage & Latency \\
\midrule
Object point-cloud construction & $\sim$100 ms \\
SAM2 preprocessing and segmentation & $\sim$200 ms \\
Semantic / observation encoding & $\sim$170 ms \\
DP3 policy inference & $\sim$155 ms \\
\bottomrule
\end{tabular}
\end{table}

Since the full perception-to-policy pipeline is slower than the low-level robot command rate, we use asynchronous execution instead of blocking robot control on every policy query. A low-frequency inference thread reads RGB-D observations, extracts object point clouds, queries the semantic field, and predicts an action chunk from the policy. In parallel, a high-frequency control thread consumes the latest available action chunk and sends smooth joint commands to the robot at a fixed servo rate. When a new chunk is not yet available, the controller holds or repeats the most recent target command as a keep-alive action. We further apply command smoothing and per-step motion limits before execution. This design decouples expensive semantic perception and policy inference from low-level command streaming, enabling continuous robot motion while using the proposed semantic field in the real-world control loop.


\section{Additional Semantic Feature Visualizations}
\label{app:additional_features}

Figure~\ref{fig:more_result} provides additional qualitative visualizations of the semantic embeddings produced by our object-centric field on different object families. Colors are obtained by projecting the queried embeddings to RGB for visualization only; consistent color patterns across instances indicate that the learned representation captures stable part-aware semantic structure.

\begin{figure*}[t]
    \centering
    \includegraphics[width=0.9\linewidth]{Images/more_result}
    \caption{Additional qualitative visualizations of queried semantic embeddings on mugs, hammers, microwaves, spoons, and teapots. Similar colors indicate similar embedding responses after RGB projection, highlighting consistent part-aware patterns across object instances.}
    \label{fig:more_result}
\end{figure*}
